\subsection{Orders and Cyclic Modules over a PID}

\begin{convention} \label{conv:standing-pid-notation} \hlt{Standing PID Notation}\\
In this and the subsequent PID-structure subsections, $R$ denotes a principal ideal domain unless otherwise stated.\\
For $r\in R$, write $\langle r\rangle=Rr$; for an $R$-module $M$ and $v\in M$, write $\langle\!\langle v\rangle\!\rangle_R=Rv$.
\end{convention}

\begin{motivation} \label{mot:pid-modules-scalar-arithmetic-controls-structure} \hlt{Scalar Arithmetic over a PID}\\
For a field $F$, $F\setminus\{0_F\}=F^\times$, so every nonzero scalar is invertible.\\
A PID need not permit universal inversion; divisibility, gcds, and lcms replace it. Module order records this scalar arithmetic.
\end{motivation}

\begin{definition} \label{def:module-order-over-pid} \hlt{Order of a Submodule or Element}\\
Let $M$ be $R$-module, $N\leq M$. An \hlt{order} of $N$ is any $\alpha\in R$ satisfying $\operatorname{Ann}_R(N)=\{r\in R:(\forall\,n\in N)\,rn=0_M\}=\langle\alpha\rangle$.\\
Write $o(N)=\alpha$; nonzero choices are understood up to associates. For $v\in M$, define $o(v)\defeq o(\langle\!\langle v\rangle\!\rangle_R)$.
\end{definition}

\begin{proposition} \label{prop:orders-direct-sum-lcm} \hlt{Orders and Direct Sums}\\
Let $n\in\N$, and let $M_1,\ldots,M_n$ be $R$-modules with $o(M_i)=\alpha_i$. Then $o(\bigoplus\limits_{i=1}^{n}M_i)=\operatorname{lcm}(\alpha_1,\ldots,\alpha_n)$.
\begin{proof}
Set $I\defeq\{1,\ldots,n\}$, $\lambda\defeq\operatorname{lcm}(\alpha_1,\ldots,\alpha_n)$. For $r\in R$, $r\in\operatorname{Ann}_R(\bigoplus\limits_{i\in I}M_i)\Longleftrightarrow(\forall\,i\in I)\,\alpha_i\mid r\Longleftrightarrow\lambda\mid r$.\\
Hence $\operatorname{Ann}_R(\bigoplus\limits_{i=1}^{n}M_i)=\langle\lambda\rangle$. Thus $o(\bigoplus\limits_{i=1}^{n}M_i)=\lambda$.\\
\hlt{Structural Commentary}: The direct-sum annihilator is the intersection of the summand annihilators.\\
\hlt{Key Insight}: Common annihilators are exactly multiples of the least common multiple.
\end{proof}
\end{proposition}

\begin{theorem} \label{thm:cyclic-modules-over-pid} \hlt{Cyclic Modules over a PID}\\
Let $M$ be an $R$-module and let $v\in M$.
\begin{enumerate}[label=(\roman*)]
\setlength{\itemsep}{0pt}
\item If $\operatorname{Ann}_R(v)=\langle\alpha\rangle$, define $\phi_v:R/\langle\alpha\rangle\to\langle\!\langle v\rangle\!\rangle_R$ by $\phi_v(r+\langle\alpha\rangle)=rv$. Then $\phi_v$ is an $R$-module isomorphism.
\item If $C=\langle\!\langle v\rangle\!\rangle_R$ and $S\leq C$, then $(\exists\,s\in R)\,S=\langle\!\langle sv\rangle\!\rangle_R$.
\item Let $o(v)=\alpha$ and $\beta\in R$. If $(\alpha,\beta)\neq(0_R,0_R)$, then $o(\beta v)=\frac{\alpha}{\gcd(\alpha,\beta)}$. If $\alpha=\beta=0_R$, then $o(\beta v)=1_R$.
\item Let $o(v)=\alpha$ and $\beta\in R$ satisfy $(\alpha,\beta)\neq(0_R,0_R)$. Then $\langle\!\langle\beta v\rangle\!\rangle_R=\langle\!\langle v\rangle\!\rangle_R\Longleftrightarrow\gcd(\alpha,\beta)\sim1_R$.
\end{enumerate}
\begin{proof}
\begin{enumerate}[label=(\roman*)]
\setlength{\itemsep}{0pt}
\item[$(i)$] Set $\psi:R\to\langle\!\langle v\rangle\!\rangle_R$, $\psi(r)=rv$. 1qThen $\operatorname{im}\psi=\langle\!\langle v\rangle\!\rangle_R$ and $\ker\psi=\operatorname{Ann}_R(v)=\langle\alpha\rangle$.\\
The induced map $\phi_v$ is an $R$-module isomorphism by Theorem \ref{thm:module-isomorphism-theorems}(i).
\item[$(ii)$] Set $\psi:R\to C$, $\psi(r)=rv$, $I\defeq\psi^{-1}(S)=\{r\in R:rv\in S\}$; $0_R\in I$ and $(\forall\,(a,b,x,y)\in R^2\times I^2)\,\psi(ax+by)=a\psi(x)+b\psi(y)\in S$.\\
Thus $I\leq{}_R R$ by Theorem \ref{thm:submodule-test}. Choose $s\in R$ with $I=\langle s\rangle$ by Definition \ref{def:principal-ideal-domain-and-ufd}(i).\\
The equality $S=\psi(I)$ holds by $S\leq\operatorname{im}\psi$, so $S=\psi(Rs)=R(sv)=\langle\!\langle sv\rangle\!\rangle_R$.
\item[$(iii)$] If $\alpha=\beta=0_R$, then $\beta v=0_M$ and $\operatorname{Ann}_R(\beta v)=R=\langle1_R\rangle$, so $o(\beta v)=1_R$. Suppose $(\alpha,\beta)\neq(0_R,0_R)$,  $d\defeq\gcd(\alpha,\beta)\neq0_R$.\\
Choose $\alpha',\beta'\in R$ with $\alpha=d\alpha'$, $\beta=d\beta'$, and $\gcd(\alpha',\beta')\sim1_R$. Choose $x,y\in R$ with $x\alpha'+y\beta'=1_R$ by Theorem \ref{thm:pid-prime-irreducible-bezout}(iii).\\
Fix $r\in R$. If $\alpha'\mid r\beta'$, then $\alpha'\mid(xr\alpha'+yr\beta')=r$; the reverse divisibility holds by multiplication by $\beta'$.\\
Thus $r\beta v=0_M\Longleftrightarrow\alpha\mid r\beta\Longleftrightarrow\alpha'\mid r$ by Proposition \ref{prop:integral-domain-cancellation} with $d\neq0_R$. Hence $\operatorname{Ann}_R(\beta v)=\langle\alpha'\rangle$ and $o(\beta v)=\alpha'=\frac{\alpha}{d}$.
\item[$(iv)$] ($\Rightarrow$) Suppose $\langle\!\langle\beta v\rangle\!\rangle_R=\langle\!\langle v\rangle\!\rangle_R$. Choose $c\in R$ with $v=c\beta v$. Then $\alpha\mid(1_R-c\beta)$, so $(\exists\,q\in R)\,q\alpha+c\beta=1_R$.\\
Hence $\gcd(\alpha,\beta)\sim1_R$ by Theorem \ref{thm:pid-prime-irreducible-bezout}(iii).
\item[$(iv)$] ($\Leftarrow$) Suppose $\gcd(\alpha,\beta)\sim1_R$. Choose $q,c\in R$ with $q\alpha+c\beta=1_R$ by Theorem \ref{thm:pid-prime-irreducible-bezout}(iii).\\
Then $v=(q\alpha+c\beta)v=c\beta v$, so $\langle\!\langle v\rangle\!\rangle_R\leq\langle\!\langle\beta v\rangle\!\rangle_R$. plaThe reverse inclusion holds by $\beta v\in\langle\!\langle v\rangle\!\rangle_R$.
\end{enumerate}
\noindent
\hlt{Structural Commentary}: Kernel quotients and preimage ideals transfer cyclic questions to ideal arithmetic.\\
\hlt{Key Insight}: Annihilators convert cyclic-module questions into divisibility in $R$.
\end{proof}
\end{theorem}

\begin{theorem} \label{thm:cyclic-composition-coprime-orders} \hlt{Composing and Decomposing Cyclic Modules}\\
Let $M$ be an $R$-module, let $n\in\N$, and set $I\defeq\{1,\ldots,n\}$.\\
Let $(\alpha_1,\ldots,\alpha_n)\in(R\setminus\{0_R\})^n$ satisfy $(\forall\,i,j\in I)\,[i\neq j\Rightarrow\gcd(\alpha_i,\alpha_j)\sim1_R]$.
\begin{enumerate}[label=(\roman*)]
\setlength{\itemsep}{0pt}
\item If $(u_1,\ldots,u_n)\in M^n$ and $(\forall\,i\in I)\,o(u_i)=\alpha_i$, then $o(\sum\limits_{i=1}^{n}u_i)\sim\prod\limits_{i=1}^{n}\alpha_i$.
\item Let $(u_1,\ldots,u_n)\in M^n$ satisfy $(\forall\,i\in I)\,o(u_i)=\alpha_i$. Then $\bigoplus\limits_{i=1}^{n}\langle\!\langle u_i\rangle\!\rangle_R=\langle\!\langle\sum\limits_{i=1}^{n}u_i\rangle\!\rangle_R$.
\item If $v\in M$ and $o(v)=\alpha=\prod\limits_{i=1}^{n}\alpha_i$, then there is a tuple $(v_1,\ldots,v_n)\in M^n$ satisfying
\begin{enumerate}[label=(\alph*)]
\setlength{\itemsep}{0pt}
\item $v=\sum\limits_{i=1}^{n}v_i$;
\item $(\forall\,i\in I)\,o(v_i)\sim\alpha_i$;
\item $\langle\!\langle v\rangle\!\rangle_R=\bigoplus\limits_{i=1}^{n}\langle\!\langle v_i\rangle\!\rangle_R$.
\end{enumerate}
\end{enumerate}
\begin{proof}
Set $\alpha\defeq\prod\limits_{i=1}^{n}\alpha_i$ and $\beta_i\defeq\frac{\alpha}{\alpha_i}=\prod\limits_{j\in I\setminus\{i\}}\alpha_j$. Then $(\forall\,i\in I)\,\gcd(\alpha_i,\beta_i)\sim1_R$ by Theorem \ref{thm:pid-is-ufd}(ii).
\begin{enumerate}[label=(\roman*)]
\setlength{\itemsep}{0pt}
\item Set $u\defeq\sum\limits_{i=1}^{n}u_i$. Then $\alpha u=0_M$, so $\alpha\in\operatorname{Ann}_R(u)$. Suppose $ru=0_M$, $i\in I$. Then $r\beta_i u_i=0_M$ by $(\forall\,j\in I\setminus\{i\})\,\alpha_j\mid\beta_i$.\\
Thus $\alpha_i\mid r\beta_i$. Choose $a_i,b_i\in R$ with $a_i\alpha_i+b_i\beta_i=1_R$ by Theorem \ref{thm:pid-prime-irreducible-bezout}(iii). Then $\alpha_i\mid r(a_i\alpha_i+b_i\beta_i)=r$.\\
Hence $(\forall\,i\in I)\,\alpha_i\mid r$. Thus $\alpha\mid r$ by Theorem \ref{thm:pid-is-ufd}(ii). Therefore $\operatorname{Ann}_R(u)=\langle\alpha\rangle$ and $o(u)\sim\alpha$.
\item Choose $((a_i,b_i))_{i=1}^{n}\in(R^2)^n$ with $(\forall\,i\in I)\,a_i\alpha_i+b_i\beta_i=1_R$ by Theorem \ref{thm:pid-prime-irreducible-bezout}(iii).\\
Suppose $\sum\limits_{i=1}^{n}x_i=0_M$ with $(\forall\,i\in I)\,x_i\in\langle\!\langle u_i\rangle\!\rangle_R$, and fix $i\in I$.\\
Then $\alpha_i x_i=0_M$ and $\beta_i x_i=0_M$, the latter by multiplying $\sum\limits_{j=1}^{n}x_j=0_M$ by $\beta_i$.\\
Hence $x_i=(a_i\alpha_i+b_i\beta_i)x_i=0_M$, and the sum is internal.\\
Set $u\defeq\sum\limits_{i=1}^{n}u_i$. The inclusion $\langle\!\langle u\rangle\!\rangle_R\leq\sum\limits_{i=1}^{n}\langle\!\langle u_i\rangle\!\rangle_R$ holds by $u$ belonging to the sum.\\
$(\forall\,i\in I)\,u_i=b_i\beta_i u$ by $\alpha_i u_i=0_M$ and $(\forall\,j\in I\setminus\{i\})\,\beta_i u_j=0_M$. Thus $\sum\limits_{i=1}^{n}\langle\!\langle u_i\rangle\!\rangle_R\leq\langle\!\langle u\rangle\!\rangle_R$.
\item If $n=1$, then $(\forall\,h\in R)\,[h\mid\beta_1\Rightarrow h\in R^\times]$ by $\beta_1=1_R$. Suppose $n\geq2$ and fix $h\in R\setminus R^\times$ satisfying $(\forall\,i\in I)\,h\mid\beta_i$.\\
Then $h\neq0_R$ by $h\mid\beta_1\neq0_R$. Choose an irreducible $p\mid h$ by Theorem \ref{thm:pid-is-ufd}(i). Then $p$ is prime by Theorem \ref{thm:pid-prime-irreducible-bezout}(ii).\\
By the prime property and $p\mid\beta_1$, choose $j\neq1$ with $p\mid\alpha_j$. By the prime property and $p\mid\beta_j$, choose $k\neq j$ with $p\mid\alpha_k$.\\
This contradicts $\gcd(\alpha_j,\alpha_k)\sim1_R$. Thus $(\forall\,h\in R)\,[(\forall\,i\in I)\,h\mid\beta_i\Rightarrow h\in R^\times]$.\\
Choose $d\in R$ with $\langle\beta_1,\ldots,\beta_n\rangle=\langle d\rangle$ by Definition \ref{def:principal-ideal-domain-and-ufd}(i).\\
The element $d$ is a common divisor of the $\beta_i$, so $d\in R^\times$ and $\langle\beta_1,\ldots,\beta_n\rangle=R$.\\
Choose $c_1,\ldots,c_n\in R$ with $\sum\limits_{i=1}^{n}c_i\beta_i=1_R$, and set $v_i\defeq c_i\beta_i v$.
\begin{enumerate}[label=(\alph*)]
\setlength{\itemsep}{0pt}
\item The identity $v=\sum\limits_{i=1}^{n}v_i$ holds by $\sum\limits_{i=1}^{n}c_i\beta_i=1_R$.
\item Fix $i\in I$. If $h\mid c_i$ and $h\mid\alpha_i$, then $(\forall\,j\in I)\,h\mid c_j\beta_j$, because $\alpha_i\mid\beta_j$ when $j\neq i$. Hence $h\mid1_R$ and $\gcd(c_i,\alpha_i)\sim1_R$.\\
Therefore $\gcd(\alpha,c_i\beta_i)\sim\beta_i$ by Theorem \ref{thm:pid-is-ufd}(ii). Thus $o(v_i)\sim\alpha_i$ by Theorem \ref{thm:cyclic-modules-over-pid}(iii).
\item The equality $\langle\!\langle v\rangle\!\rangle_R=\bigoplus\limits_{i=1}^{n}\langle\!\langle v_i\rangle\!\rangle_R$ holds by (ii), $v=\sum\limits_{i=1}^{n}v_i$, and $(\forall\,i\in I)\,o(v_i)\sim\alpha_i$.
\end{enumerate}
\end{enumerate}
\noindent
\hlt{Structural Commentary}: Bézout projections compose the summands; complementary products construct the reverse decomposition.\\
\hlt{Key Insight}: Coprime orders isolate each cyclic component from the total sum.
\end{proof}
\end{theorem}

\begin{proposition} \label{prop:prime-divisor-prime-order-submodule} \hlt{Prime-Order Cyclic Submodules from Prime Divisors}\\
Let $M$ be a finitely generated torsion $R$-module with $o(M)=\mu$. If $p\mid\mu$ is prime, then $(\exists\,w\in M)\,[o(w)\sim p\land Rw\cong R/\langle p\rangle]$.
\begin{proof}
$\operatorname{Ann}_R(M)\neq\{0_R\}$ by Proposition \ref{prop:finitely-generated-torsion-module-annihilator}, so $\mu\neq0_R$. Choose $q\in R$ with $\mu=pq$. If $q\in\langle\mu\rangle$, choose $a\in R$ with $q=\mu a=pqa$.\\
Then $1_R=pa$ by Proposition \ref{prop:integral-domain-cancellation} with $q\neq0_R$, contrary to $p\notin R^\times$. Thus $q\notin\operatorname{Ann}_R(M)$.\\
Choose $v\in M$ with $qv\neq0_M$ by the definition of $\operatorname{Ann}_R(M)$, and set $w\defeq qv$. Then $pw=\mu v=0_M$, so $o(w)\mid p$.\\
The element $o(w)$ is a nonunit by $w\neq0_M$. Write $p=o(w)c$. Then $p\mid o(w)$ or $p\mid c$ by primality of $p$.\\
If $p\mid c$, write $c=pe$; then $1_R=o(w)e$ by Proposition \ref{prop:integral-domain-cancellation} in $p=o(w)pe$, contrary to $o(w)\notin R^\times$.\\
Thus $p\mid o(w)$, and $o(w)\sim p$ by Proposition \ref{prop:principal-ideals-divisibility-associates}(ii)--(iii). Finally, $Rw\cong R/\langle p\rangle$ by Theorem \ref{thm:cyclic-modules-over-pid}(i).\\
\hlt{Structural Commentary}: A complementary factor constructs a nonzero element annihilated by the selected prime.\\
\hlt{Key Insight}: The complementary factor is not a global annihilator.
\end{proof}
\end{proposition}
